\documentclass[../../main.tex]{subfiles}
\begin{document}
% SOURCE p=211; printed=196; object=ch06-inequalities-exercises-title; region=upper
\Needspace{12\baselineskip}
\subsection*{习题}
% SOURCE p=211; printed=196; object=ch06-ex26; number=26; region=upper
\begin{exercise}\label{exercise:ch06-kernel-endpoints}
补全\cref{theorem:ch06-kernel-bound}在 $p=1$ 和 $p=\infty$ 时的证明。
\end{exercise}
% SUPPLEMENT: id=v14-ch06-ex26-solution; parent=ch06-ex26; kind=solution+motivation+hypotheses; origin=independent-derivation
\begin{solution}
\textbf{两个端点的完整证明（重构新增）。}令 $H(x)=\int_Y|K(x,y)f(y)|\,\mathrm d\nu(y)$。当 $p=1$ 时，非负函数的 Tonelli 定理给出
\begin{align*}
\int_XH\,\mathrm d\mu
=\int_Y|f(y)|\left(\int_X|K(x,y)|\,\mathrm d\mu(x)\right)\mathrm d\nu(y)
\leqslant C\|f\|_1<\infty.
\end{align*}
因此 $H$ 几乎处处有限，$Tf$ 几乎处处绝对收敛；又 $|Tf|\leqslant H$，故 $\|Tf\|_1\leqslant C\|f\|_1$。

当 $p=\infty$ 时，可取处处满足 $|f|\leqslant\|f\|_\infty$ 的可测代表。对行积分界成立的每个 $x$，
\begin{align*}
H(x)\leqslant\|f\|_\infty\int_Y|K(x,y)|\,\mathrm d\nu(y)
\leqslant C\|f\|_\infty<\infty.
\end{align*}
所以 $Tf$ 几乎处处绝对收敛，且 $\|Tf\|_\infty\leqslant C\|f\|_\infty$。两种情形的可测性由非负截面积分得到；在绝对积分无穷的可测零集上令 $Tf=0$。若换回任意等价代表，改变只在固定的 $\nu$ 零集上，不影响各个截面积分。绝对收敛还保证线性。
\end{solution}
\begin{remark}
\textbf{动机与条件说明（重构新增）。}$L^1$ 估计把输出变量先积分掉，只用列积分界；$L^\infty$ 估计先取输入函数的统一界，只用行积分界。两者都直接控制绝对积分，因此同时解决存在性与有界性。零函数和零测度空间包含在估计中，无需单位范数归一化。联合可测性及 $\sigma$ 有限性沿用原定理以适用 Tonelli；在 $L^\infty$ 的单个截面估计中，并未另行使用列积分条件。
\end{remark}
% END SUPPLEMENT: id=v14-ch06-ex26-solution
% SOURCE p=211; printed=196; object=ch06-ex27; number=27; region=upper
\begin{exercise}[（Hilbert 不等式）]\label{exercise:ch06-hilbert-inequality}
算子 $Tf(x)=\int_0^\infty(x+y)^{-1}f(y)\,\mathrm dy$ 对 $1<p<\infty$ 满足 $\|Tf\|_p\leqslant C_p\|f\|_p$，其中
\begin{align*}
C_p=\int_0^\infty\frac{x^{-1/p}}{x+1}\,\mathrm dx.
\end{align*}
（对熟悉围道积分的读者：证明 $C_p=\pi\csc(\pi/p)$。）
\end{exercise}
% SUPPLEMENT: id=v14-ch06-ex27-solution; parent=ch06-ex27; kind=solution+motivation+hypotheses; origin=independent-derivation
\begin{solution}
\textbf{算子估计与常数计算（重构新增）。}核 $K(x,y)=(x+y)^{-1}$ 可测、对称，且满足负一次齐次性。在 $(0,1)$ 上，其加权剖面不超过 $x^{-1/p}$，在 $(1,\infty)$ 上不超过 $x^{-1-1/p}$，二者分别因 $p>1$ 与 $p<\infty$ 可积。故\cref{theorem:ch06-homogeneous-kernel}给出绝对收敛几乎处处及所求估计。

下面完成题中附带的复分析问，使用留数定理。置 $a=1-1/p\in(0,1)$，在沿正实轴割开的平面取 $0<\arg z<2\pi$ 的对数分支，考虑
\begin{align*}
F(z)=\frac{z^{a-1}}{1+z}.
\end{align*}
取内半径 $0<\rho<1/2$、外半径 $R>2$ 的正向钥匙孔围道：沿正轴上侧由 $\rho$ 到 $R$，外圆逆时针，再沿正轴下侧回到 $\rho$，最后内圆顺时针闭合。严格地可先把两条射线取在角度 $\theta$ 和 $2\pi-\theta$，再令 $\theta\downarrow0$；固定 $\rho,R$ 上的连续性与有界性允许这个极限。围道内唯一极点为 $-1$，留数为 $e^{i\pi(a-1)}$。两条直线段之和为
\begin{align*}
\bigl(1-e^{2\pi ia}\bigr)\int_\rho^R\frac{x^{a-1}}{1+x}\,\mathrm dx.
\end{align*}
外圆积分模至多 $2\pi R^a/(R-1)$，随 $R\to\infty$ 趋零，因为 $a<1$；内圆积分模至多 $2\pi\rho^a/(1-\rho)$，随 $\rho\downarrow0$ 趋零，因为 $a>0$。留数定理于是给出
\begin{align*}
\bigl(1-e^{2\pi ia}\bigr)\int_0^\infty\frac{x^{a-1}}{1+x}\,\mathrm dx
=2\pi i e^{i\pi(a-1)}.
\end{align*}
由 $1-e^{2\pi ia}=-2i e^{i\pi a}\sin(\pi a)$ 及 $e^{i\pi(a-1)}=-e^{i\pi a}$，得到
\begin{align*}
C_p=\frac{\pi}{\sin(\pi a)}=\frac{\pi}{\sin(\pi/p)}.
\end{align*}
\end{solution}
\begin{remark}
\textbf{动机与条件说明（重构新增）。}分母 $x+y$ 的齐次性把两变量问题化为一个加权积分，正是前一定理的适用结构。两个严格端点分别控制零附近与无穷远的可积性；$p=1$ 或 $\infty$ 时这个常数发散，不能直接套用估计。围道法中的两个小量也分别对应 $0<a<1$ 的两端限制。复分析部分以留数定理为已知前提，符合原题明确给出的附加背景；算子有界性的证明不依赖它，也不依赖先知常数的闭式。
\end{remark}
% END SUPPLEMENT: id=v14-ch06-ex27-solution
% SOURCE p=211; printed=196; object=ch06-ex28; number=28; region=middle
\begin{exercise}\label{exercise:ch06-fractional-hardy}
令 $I_\alpha$ 为\cref{exercise:ch02-fractional-integral}中的 $\alpha$ 阶分数积分算子，并令 $J_\alpha f(x)=x^{-\alpha}I_\alpha f(x)$。
\begin{enumerate}
\item\label{item:ch06-fractional-hardy-bound} 当 $1<p\leqslant\infty$ 时，$J_\alpha$ 在 $L^p((0,\infty))$ 上有界；更精确地，
\begin{align*}
\|J_\alpha f\|_p\leqslant\frac{\Gamma(1-p^{-1})}{\Gamma(\alpha+1-p^{-1})}\|f\|_p.
\end{align*}
\item\label{item:ch06-fractional-hardy-endpoint} 存在 $f\in L^1((0,\infty))$，使 $J_1f\notin L^1((0,\infty))$。
\end{enumerate}
\end{exercise}
% SUPPLEMENT: id=v14-ch06-ex28-solution; parent=ch06-ex28; kind=solution+motivation+hypotheses; origin=independent-derivation
\begin{solution}
\textbf{分数积分核与端点反例（重构新增）。}沿用定义中的 $\alpha>0$，对 $0<t<x$ 取
\begin{align*}
K(t,x)=\frac{x^{-\alpha}(x-t)^{\alpha-1}}{\Gamma(\alpha)},
\end{align*}
在其余点令 $K=0$，尤其不在对角线上计算可能出现的 $0^{\alpha-1}$。这是可测的负一次齐次核，且
\begin{align*}
\int_0^\infty|K(t,1)|t^{-1/p}\,\mathrm dt
&=\frac1{\Gamma(\alpha)}\int_0^1(1-t)^{\alpha-1}t^{-1/p}\,\mathrm dt\\
&=\frac{\Gamma(1-1/p)}{\Gamma(\alpha+1-1/p)}.
\end{align*}
最后一步为\cref{exercise:ch02-beta-integral}的公式。为明确其依据，对任意 $b,\alpha>0$，Tonelli 与换元 $u=rt,v=r(1-t)$、Jacobian 为 $r$，给出
\begin{align*}
\Gamma(b)\Gamma(\alpha)
&=\int_0^\infty\int_0^\infty e^{-(u+v)}u^{b-1}v^{\alpha-1}\,\mathrm du\,\mathrm dv\\
&=\left(\int_0^\infty e^{-r}r^{b+\alpha-1}\,\mathrm dr\right)
\left(\int_0^1t^{b-1}(1-t)^{\alpha-1}\,\mathrm dt\right).
\end{align*}
取 $b=1-1/p>0$ 即得该式；$p=\infty$ 时 $b=1$ 也成立。由\cref{theorem:ch06-homogeneous-kernel}，积分对一般 $L^p$ 函数绝对收敛几乎处处，定义了 $J_\alpha$，并满足 (a) 的界。

对 (b)，取 $f=\chi_{(0,1)}$，则 $\|f\|_1=1$、$\Gamma(1)=1$，并且
\begin{align*}
J_1f(x)=\begin{cases}1,&0<x\leqslant1,\\x^{-1},&x>1.\end{cases}
\end{align*}
其尾积分发散，故不属于 $L^1$。
\end{solution}
\begin{remark}
\textbf{动机与条件说明（重构新增）。}因子 $x^{-\alpha}$ 抵消分数积分的尺度，使核成为负一次齐次。$\alpha>0$ 控制积分上端的 $(1-t)^{\alpha-1}$ 奇性，$p>1$ 控制下端的 $t^{-1/p}$ 奇性，Gamma 比值由这两种可积性共同决定。原第二章先对连续函数定义 $I_\alpha$；这里通过绝对收敛的积分式扩展到 $L^p$ 等价类，不暗中假设 $f$ 连续。第二问说明 $p=1$ 的有界性不能普遍补入；$p=\infty$ 则完全包含在证明中。
\end{remark}
% END SUPPLEMENT: id=v14-ch06-ex28-solution
% SOURCE p=211; printed=196; object=ch06-ex29; number=29; region=middle
\begin{exercise}\label{exercise:ch06-weighted-hardy}
设 $1\leqslant p<\infty$，$r>0$，且 $h$ 为 $(0,\infty)$ 上的非负可测函数。则
\begin{align*}
\int_0^\infty x^{-r-1}\left(\int_0^x h(y)\,\mathrm dy\right)^p\mathrm dx
&\leqslant\left(\frac pr\right)^p\int_0^\infty x^{p-r-1}h(x)^p\,\mathrm dx,\\
\int_0^\infty x^{r-1}\left(\int_x^\infty h(y)\,\mathrm dy\right)^p\mathrm dx
&\leqslant\left(\frac pr\right)^p\int_0^\infty x^{p+r-1}h(x)^p\,\mathrm dx.
\end{align*}
（在\cref{theorem:ch06-homogeneous-kernel}中取 $K(x,y)=x^{\beta-1}y^{-\beta}\chi_{(0,\infty)}(y-x)$、$f(x)=x^\gamma h(x)$、$g(x)=x^\delta h(x)$，选择适当的 $\beta,\gamma,\delta$。）
\end{exercise}
% SUPPLEMENT: id=v14-ch06-ex29-solution; parent=ch06-ex29; kind=solution+motivation+hypotheses; origin=independent-derivation
\begin{solution}
\textbf{两次权重配平（重构新增）。}每一式的右端若无穷，该式在扩展实数意义下成立；以下分别假设相应右端有限。记
\begin{align*}
K_\beta(x,y)=x^{\beta-1}y^{-\beta}\chi_{\{x<y\}}.
\end{align*}
它可测且负一次齐次。对任意指数 $P\in[1,\infty]$，只要 $\beta>1/P$，就有
\begin{align*}
C_{\beta,P}=\int_0^1x^{\beta-1-1/P}\,\mathrm dx=\frac1{\beta-1/P}.
\end{align*}
为证明第一式，取 $P=p$、$\beta=(r+1)/p$、$\gamma=1-\beta$，令 $f(x)=x^{1-\beta}h(x)$。则
\begin{align*}
T_\beta f(y)=y^{-\beta}\int_0^y h(x)\,\mathrm dx,\qquad
\|f\|_p^p=\int_0^\infty x^{p-r-1}h(x)^p\,\mathrm dx.
\end{align*}
此时 $\beta-1/p=r/p>0$，故 $\|T_\beta f\|_p\leqslant(p/r)\|f\|_p$；取 $p$ 次幂，注意 $-p\beta=-r-1$，即得第一式。

为证明第二式，让 $P$ 为 $p$ 的共轭指数，即 $1/P=1-1/p$；特别地，$p=1$ 时 $P=\infty$。取 $\beta=1+(r-1)/p$、$\delta=\beta$，令 $g(y)=y^\beta h(y)$。现在应用定理中的转置算子：
\begin{align*}
S_\beta g(x)=x^{\beta-1}\int_x^\infty h(y)\,\mathrm dy,\qquad
\|g\|_p^p=\int_0^\infty y^{p+r-1}h(y)^p\,\mathrm dy.
\end{align*}
因为 $\beta-1/P=r/p>0$，同一定理给 $\|S_\beta g\|_p\leqslant(p/r)\|g\|_p$。取 $p$ 次幂并用 $p(\beta-1)=r-1$，得到第二式。每次所用输入确属于相应 $L^p$，正是该式右端有限的假设；定理同时保证需要的内积分几乎处处有限。
\end{solution}
\begin{remark}
\textbf{动机与条件说明（重构新增）。}这些幂权的作用是把非加权 Hardy 核调整到所需尺度。两个不等式分别使用正向和转置算子，因而 $\beta$ 不能机械地取成同一个数。$r>0$ 等价于关键剖面的可积条件，并使常数 $p/r$ 有限；$1\leqslant p<\infty$ 允许使用范数定理再取有限 $p$ 次幂，其中 $p=1$ 通过转置一侧的 $P=\infty$ 合法处理。$h\geqslant0$ 使原式即便内积分无穷也有明确含义；右端有限时才需要实际算子论证。可测性保证幂权乘积和积分均可测。
\end{remark}
% END SUPPLEMENT: id=v14-ch06-ex29-solution
% SOURCE p=211; printed=196; object=ch06-ex30; number=30; region=lower
\begin{exercise}\label{exercise:ch06-mellin-laplace-bound}
设 $K$ 为 $(0,\infty)$ 上的非负可测函数，且对 $0<s<1$ 有
\begin{align*}
\phi(s)=\int_0^\infty K(x)x^{s-1}\,\mathrm dx<\infty.
\end{align*}
\begin{enumerate}
\item\label{item:ch06-mellin-bilinear} 若 $1<p<\infty$、$p^{-1}+q^{-1}=1$，且 $f,g$ 为 $(0,\infty)$ 上的非负可测函数，则（各积分均在 $(0,\infty)$ 上）
\begin{align*}
\iint K(xy)f(x)g(y)\,\mathrm dx\,\mathrm dy
\leqslant\phi(p^{-1})\left(\int x^{p-2}f(x)^p\,\mathrm dx\right)^{1/p}
\left(\int g(x)^q\,\mathrm dx\right)^{1/q}.
\end{align*}
\item\label{item:ch06-laplace-ltwo} 算子 $Tf(x)=\int_0^\infty K(xy)f(y)\,\mathrm dy$ 在 $L^2((0,\infty))$ 上有界，范数不超过 $\phi(1/2)$。（一个有趣的特例：若 $K(x)=e^{-x}$，则 $T$ 为 Laplace 变换，且 $\phi(s)=\Gamma(s)$。）
\end{enumerate}
\end{exercise}
% SUPPLEMENT: id=v14-ch06-ex30-solution; parent=ch06-ex30; kind=solution+motivation+hypotheses; origin=independent-derivation
\begin{solution}
\textbf{倒数换元后的齐次核（重构新增）。}先假设 (a) 右边两个范数均有限。令
\begin{align*}
L(t,y)=t^{-1}K(y/t),\qquad F(t)=t^{-1}f(1/t).
\end{align*}
$L$ 为非负可测的负一次齐次核；Lebesgue 可测情形的联合可测性由局部光滑坐标变换及完成形式的 Fubini 保证。换元 $u=1/t$ 得
\begin{align*}
\|F\|_p^p&=\int_0^\infty u^{p-2}f(u)^p\,\mathrm du,\\
\int_0^\infty L(t,1)t^{-1/p}\,\mathrm dt
&=\int_0^\infty K(u)u^{1/p-1}\,\mathrm du=\phi(1/p).
\end{align*}
由\cref{theorem:ch06-homogeneous-kernel}，函数
\begin{align*}
AF(y)=\int_0^\infty L(t,y)F(t)\,\mathrm dt
=\int_0^\infty K(xy)f(x)\,\mathrm dx
\end{align*}
属于 $L^p$，且 $\|AF\|_p\leqslant\phi(1/p)\|F\|_p$。最后对 $AF$ 与 $g$ 用 Hölder，即得 (a)。对一般非负可测 $f,g$，取
\begin{align*}
f_n=\min(f,n)\chi_{[1/n,n]},\qquad g_n=\min(g,n)\chi_{[1/n,n]}.
\end{align*}
它们单调增加到原函数，相关范数有限。对截断函数应用已证不等式再用单调收敛，得到扩展实数意义的结论；若一项范数为零，则相应函数几乎处处为零，左边为零，单独处理而不进行未定义的数值乘法。

对 (b)，令 $Uf(t)=t^{-1}f(1/t)$。同一换元给 $\|Uf\|_2=\|f\|_2$，且 $U^2f=f$。把上面的齐次核定理取 $p=2$，便得 $A$ 在 $L^2$ 上有界，范数不超过 $\phi(1/2)$；而 $Tf=A(Uf)$。定理对 $|Uf|$ 的估计保证所写积分几乎处处绝对收敛，因此这个恒等式对复值 $L^2$ 函数也成立。故 $\|Tf\|_2\leqslant\phi(1/2)\|f\|_2$，且 $T$ 线性、与代表选择无关。若 $K(x)=e^{-x}$，由 Gamma 函数定义，$\phi(s)=\Gamma(s)$，从而还得到 $\|T\|\leqslant\Gamma(1/2)=\sqrt\pi$。
\end{solution}
\begin{remark}
\textbf{动机与条件说明（重构新增）。}乘积型核 $K(xy)$ 本身通常不齐次；倒数换元把乘积变为比值，额外的 $t^{-1}$ 则赋予恰好的齐次次数。加权 $L^p$ 范数因此自然出现，$p=2$ 时权消失，成为等距变换。$1<p<\infty$ 使 $1/p\in(0,1)$，可用假设中的有限 $\phi(1/p)$；其实每个固定指数只需这一点的积分有限，原题在整个区间上的假设一次覆盖所有指数。非负性允许 (a) 的 Tonelli 与单调截断；(b) 的输入可复值，因为先控制绝对值。若 $\phi(1/p)=0$，正权积分为零说明 $K=0$ 几乎处处，对每个固定 $y>0$ 的伸缩也保持这个性质，因此所有相关积分为零，退化情形同样成立。
\end{remark}
% END SUPPLEMENT: id=v14-ch06-ex30-solution
% SOURCE p=211; printed=196; object=ch06-ex31; number=31; region=bottom
\begin{exercise}[（广义 Hölder 不等式）]\label{exercise:ch06-generalized-holder}
设 $1\leqslant p_j\leqslant\infty$ 且 $\sum_{j=1}^n p_j^{-1}=r^{-1}\leqslant1$。若 $f_j\in L^{p_j}$，$j=1,\ldots,n$，则 $\prod_{j=1}^nf_j\in L^r$，且
\begin{align*}
\left\|\prod_{j=1}^nf_j\right\|_r\leqslant\prod_{j=1}^n\|f_j\|_{p_j}.
\end{align*}
（先证明 $n=2$ 的情形。）
\end{exercise}
% SUPPLEMENT: id=v14-ch06-ex31-solution; parent=ch06-ex31; kind=solution+motivation+hypotheses; origin=independent-derivation
\begin{solution}
\textbf{幂次转换与归纳（重构新增）。}先设 $n=2$ 且 $r<\infty$。由 $1/p_1+1/p_2=1/r$，数 $p_1/r,p_2/r$ 属于 $[1,\infty]$ 且互为共轭指数；当某个 $p_j=\infty$ 时，另一个比值为一。对 $|f_1|^r,|f_2|^r$ 用 Hölder，得到
\begin{align*}
\int|f_1f_2|^r\leqslant
\bigl\||f_1|^r\bigr\|_{p_1/r}\bigl\||f_2|^r\bigr\|_{p_2/r}
=\|f_1\|_{p_1}^r\|f_2\|_{p_2}^r.
\end{align*}
有限 $p_j$ 的最后等式按积分计算，无穷 $p_j$ 则用本质上确界的幂次。取 $r$ 次方根即可。若 $r=\infty$，两个 $p_j$ 都是无穷；合并两个例外零集，逐点乘上界即得无穷范数估计。

对一般 $n$，令 $1/r_k=\sum_{j=1}^k1/p_j$，则每个 $r_k\in[1,\infty]$，且 $1/r_k=1/r_{k-1}+1/p_k$。从 $k=2$ 起逐次对前 $k-1$ 项的乘积和第 $k$ 项使用二因子结论，得到
\begin{align*}
\left\|\prod_{j=1}^kf_j\right\|_{r_k}\leqslant\prod_{j=1}^k\|f_j\|_{p_j}.
\end{align*}
取 $k=n$ 即得结果。$n=1$ 时就是恒等式。证明没有除以任何范数，所以某个因子为零几乎处处的情形也已包含。
\end{solution}
\begin{remark}
\textbf{动机与条件说明（重构新增）。}指数倒数相加反映乘积的可积性损失；对 $r$ 次幂使用普通 Hölder 正好把这条关系变成共轭条件。$\sum1/p_j\leqslant1$ 保证所有归纳中间指数至少为一，结果仍处在本节的赋范 $L^p$ 范围。可测性使有限乘积可测，因子范数有限使右边有限；无需任何有限测度或 $\sigma$ 有限假设。所有指数无穷时要直接用本质上确界，不能把有限积分证明形式化地令 $r=\infty$。
\end{remark}
% END SUPPLEMENT: id=v14-ch06-ex31-solution
% SOURCE p=212; printed=197; object=ch06-ex32; number=32; region=top
\begin{exercise}\label{exercise:ch06-square-integrable-kernel}
设 $(X,\mathscr M,\mu)$、$(Y,\mathscr N,\nu)$ 是 $\sigma$ 有限测度空间，且 $K\in L^2(\mu\times\nu)$。若 $f\in L^2(\nu)$，则积分
\begin{align*}
Tf(x)=\int_YK(x,y)f(y)\,\mathrm d\nu(y)
\end{align*}
对几乎所有 $x\in X$ 绝对收敛；而且 $Tf\in L^2(\mu)$，$\|Tf\|_2\leqslant\|K\|_2\|f\|_2$。
\end{exercise}
% SUPPLEMENT: id=v14-ch06-ex32-solution; parent=ch06-ex32; kind=solution+motivation+hypotheses; origin=independent-derivation
\begin{solution}
\textbf{截面的 Schwarz 估计（重构新增）。}由 Tonelli 定理，
\begin{align*}
\int_X\left(\int_Y|K(x,y)|^2\,\mathrm d\nu(y)\right)\mathrm d\mu(x)=\|K\|_2^2<\infty.
\end{align*}
故对几乎所有 $x$，$K(x,\cdot)\in L^2(\nu)$。在这些截面上使用 Schwarz 不等式，得到
\begin{align*}
H(x):=\int_Y|K(x,y)f(y)|\,\mathrm d\nu(y)
\leqslant\left(\int_Y|K(x,y)|^2\,\mathrm d\nu(y)\right)^{1/2}\|f\|_2<\infty.
\end{align*}
因而 $Tf$ 几乎处处绝对收敛，并且
\begin{align*}
\int_X|Tf(x)|^2\,\mathrm d\mu(x)\leqslant\int_XH(x)^2\,\mathrm d\mu(x)
\leqslant\|K\|_2^2\|f\|_2^2.
\end{align*}
积分的可测性由非负实部、虚部的截面积分得到，在例外零集上补零。若 $K$ 换成相同乘积几乎处处类的另一个代表，Tonelli 保证几乎所有截面也相同几乎处处；若换 $f$ 的代表，则每个合格截面积分不变。因此算子定义于所述等价类，取平方根得到范数估计。
\end{solution}
\begin{remark}
\textbf{动机与条件说明（重构新增）。}平方可积核让每个截面都成为可与 $L^2$ 输入配对的权重，而所有截面范数的平方又可积。两次平方结构相配，正是常数 $\|K\|_2$ 的来源。$\sigma$ 有限性及联合可测性保证 Tonelli 适用；$K,f$ 的平方可积性分别控制核与输入。没有要求 $K$ 具备前面行列 $L^1$ 的统一界，也没有声称 $Kf$ 在整个乘积空间属于 $L^1$；所需结论只是截面积分绝对收敛与输出的 $L^2$ 控制。
\end{remark}
% END SUPPLEMENT: id=v14-ch06-ex32-solution
% SOURCE p=212; printed=197; object=ch06-ex33; number=33; region=upper
\begin{exercise}\label{exercise:ch06-weighted-primitive-czero}
给定 $1<p<\infty$，令 $Tf(x)=x^{-1/p}\int_0^x f(t)\,\mathrm dt$。若 $p^{-1}+q^{-1}=1$，则 $T$ 是从 $L^q((0,\infty))$ 到 $C_0((0,\infty))$ 的有界线性映射。
\end{exercise}
% SUPPLEMENT: id=v14-ch06-ex33-solution; parent=ch06-ex33; kind=solution+motivation+hypotheses; origin=independent-derivation
\begin{solution}
\textbf{连续性与两端消失（重构新增）。}Hölder 不等式给 $f\in L^1(0,A)$ 对每个有限 $A>0$ 成立，且
\begin{align*}
|Tf(x)|\leqslant x^{-1/p}x^{1/p}\|f\chi_{(0,x)}\|_q
=\|f\chi_{(0,x)}\|_q\leqslant\|f\|_q.
\end{align*}
原函数 $F(x)=\int_0^x f$ 在每个有限闭区间上绝对连续，故 $x^{-1/p}F(x)$ 在正半轴连续。由于 $q<\infty$，$\int_0^x|f|^q\to0$，上式给 $Tf(x)\to0$ 当 $x\downarrow0$。

对另一端，固定 $A>0$。当 $x>A$ 时，拆分积分并再次用 Hölder，得到
\begin{align*}
|Tf(x)|
&\leqslant x^{-1/p}\left|\int_0^A f(t)\,\mathrm dt\right|
+x^{-1/p}(x-A)^{1/p}\|f\chi_{(A,\infty)}\|_q\\
&\leqslant x^{-1/p}\int_0^A|f(t)|\,\mathrm dt+\|f\chi_{(A,\infty)}\|_q.
\end{align*}
因 $p<\infty$，令 $x\to\infty$ 后第一项趋零；再令 $A\to\infty$，因 $q<\infty$，尾范数趋零。因此 $Tf$ 在无穷远也趋零。在正半轴上，这两端消失连同连续性意味着 $Tf\in C_0$：每个正水平集闭且被夹在某个紧区间 $[\delta,A]\subset(0,\infty)$ 内。积分与代表无关且线性，第一条估计给 $\|T\|\leqslant1$。
\end{solution}
\begin{remark}
\textbf{动机与条件说明（重构新增）。}因子 $x^{-1/p}$ 正好抵消区间示性函数的 $L^p$ 范数，产生统一上界。严格内部指数使 $p,q$ 都有限：$q<\infty$ 用于小区间与尾部的范数绝对连续性，$p<\infty$ 用于固定初段的衰减。端点确会失败：$p=1,q=\infty$ 时取 $f=1$ 得 $Tf=1$；$p=\infty,q=1$ 时取 $f=\chi_{(0,1)}$，输出在无穷远趋一。$C_0((0,\infty))$ 要求离开所有紧集时趋零，故不仅要查 $x\to\infty$，也须查 $x\downarrow0$。若需核对常数，取 $f=a^{-1/q}\chi_{(0,a)}$，在 $x=a$ 达到 $|Tf(a)|=1=\|f\|_q$，故算子范数实际为一。
\end{remark}
% END SUPPLEMENT: id=v14-ch06-ex33-solution
% SOURCE p=212; printed=197; object=ch06-ex34; number=34; region=upper-middle
\begin{exercise}\label{exercise:ch06-weighted-derivative-boundary}
若对每个 $0<\varepsilon<1$，$f$ 在 $[\varepsilon,1]$ 上绝对连续，且 $\int_0^1x|f'(x)|^p\,\mathrm dx<\infty$，则：$p>2$ 时，$\lim_{x\to0}f(x)$ 存在且有限；$p=2$ 时，$|f(x)|/|\log x|^{1/2}\to0$；$p<2$ 时，$|f(x)|/x^{1-2/p}\to0$。
\end{exercise}
% SUPPLEMENT: id=v14-ch06-ex34-solution; parent=ch06-ex34; kind=solution+motivation+hypotheses; origin=independent-derivation
\begin{solution}
\textbf{三种增长率及 $p=1$ 端点（重构新增）。}按本节赋范 $L^p$ 的背景，在这里取 $1\leqslant p<\infty$；原题未重写这个下限，范围的必要澄清见随后的注。局部绝对连续性使任意 $0<x<\delta\leqslant1$ 上都有 $f(\delta)-f(x)=\int_x^\delta f'(t)\,\mathrm dt$。先设 $p>1$，$q=p/(p-1)$，记
\begin{align*}
A(\delta)=\left(\int_0^\delta t|f'(t)|^p\,\mathrm dt\right)^{1/p},
\qquad A(\delta)\longrightarrow0\quad(\delta\downarrow0).
\end{align*}
加权 Hölder 给
\begin{align*}
|f(x)-f(\delta)|
\leqslant\int_x^\delta t^{1/p}|f'(t)|t^{-1/p}\,\mathrm dt
\leqslant A(\delta)\left(\int_x^\delta t^{-1/(p-1)}\,\mathrm dt\right)^{1/q}.
\end{align*}
若 $p>2$，令 $a=(p-2)/(p-1)>0$。对 $0<x<y<\delta$，同一估计给
\begin{align*}
|f(x)-f(y)|\leqslant A(\delta)\left(\frac{\delta^a}{a}\right)^{1/q}\longrightarrow0
\quad(\delta\downarrow0).
\end{align*}
因此 $f(x)$ 在趋零时满足 Cauchy 准则；由实数或复数的完备性，极限存在且有限。

若 $p=2$，上式成为
\begin{align*}
|f(x)|\leqslant|f(\delta)|+A(\delta)\bigl(\log(\delta/x)\bigr)^{1/2}.
\end{align*}
固定 $\delta\in(0,1)$，除以 $|\log x|^{1/2}$ 后令 $x\downarrow0$，得到比值的上极限至多 $A(\delta)$；再令 $\delta\downarrow0$，上极限为零，所求比值非负，故极限为零。

若 $1<p<2$，令 $b=(2-p)/(p-1)>0$，则
\begin{align*}
\int_x^\delta t^{-1/(p-1)}\,\mathrm dt
=\frac{x^{-b}-\delta^{-b}}b\leqslant b^{-1}x^{-b}.
\end{align*}
因 $-b/q=1-2/p$，有
\begin{align*}
\frac{|f(x)|}{x^{1-2/p}}
\leqslant |f(\delta)|x^{2/p-1}+b^{-1/q}A(\delta).
\end{align*}
先 $x\downarrow0$ 再 $\delta\downarrow0$，得到所需极限。

最后，当 $p=1$ 时不能使用上面的有限共轭指数计算；直接估计
\begin{align*}
x|f(x)|\leqslant x|f(\delta)|+\int_x^\delta x|f'(t)|\,\mathrm dt
\leqslant x|f(\delta)|+\int_0^\delta t|f'(t)|\,\mathrm dt.
\end{align*}
依次令 $x\downarrow0$、$\delta\downarrow0$，便得 $x|f(x)|\to0$，恰为 $p=1$ 时的第三个结论。
\end{solution}
\begin{remark}
\textbf{动机与逐条件说明（重构新增）。}局部绝对连续性把函数增量精确还原为导数积分；仅有几乎处处可微不能保证这个等式。带权导数积分控制 $t^{1/p}f'$，于是与之配对的权 $t^{-1/p}$ 的共轭次幂在零附近的可积性决定临界值 $p=2$：可积时得有限边界值，临界时出现对数，低于临界时出现幂增长。积分有限使 $A(\delta)\to0$，这是结论为小 $o$ 而非仅大 $O$ 的原因。固定 $\delta$ 后才令 $x\to0$，最后再令 $\delta\to0$，不能把两个极限顺序省掉。

原题的 $p<2$ 在这里按 $1\leqslant p<2$ 解释，不能无条件推广至 $0<p<1$。确实，固定这样的 $p$，令 $x_n=2^{-n}$、$M_n=2^{n^2}$，选择
\begin{align*}
0<d_n\leqslant\min\left\{\frac{x_n}4,
\left(\frac{2^{-n}}{4x_nM_n^p}\right)^{1/(1-p)}\right\}.
\end{align*}
在互不相交的区间 $(x_n-d_n,x_n+d_n)$ 上作高度 $M_n$、中心为 $x_n$、端点为零的三角形函数，其余处为零，记总函数为 $f$。每个 $[\varepsilon,1]$ 只遇到有限个三角形，所以 $f$ 在该区间绝对连续。每个三角形内部的 $|f'|=M_n/d_n$ 几乎处处，且 $t\leqslant2x_n$，因而
\begin{align*}
\int_0^1t|f'(t)|^p\,\mathrm dt
\leqslant\sum_n4x_nM_n^pd_n^{1-p}\leqslant\sum_n2^{-n}<\infty.
\end{align*}
但 $f(x_n)/x_n^{1-2/p}=2^{n^2-n(2/p-1)}\to\infty$。这说明必须交代赋范指数范围，而不能把原题未重写的范围默认为所有正数。
\end{remark}
% END SUPPLEMENT: id=v14-ch06-ex34-solution
\end{document}
